%%%%%%%%%%%%%%%%%%%%%%%%%%%%%%%%%%%%%%%%%
% Libro de investigación — plantilla en español con normas APA 7
%
% Diseño basado en "The Legrand Orange Book" (LaTeXTemplates.com),
% Mathias Legrand y Vel. Licencia CC BY-NC-SA 3.0.
%
% Compilación en Overleaf: pdfLaTeX (por defecto). Overleaf ejecuta biber
% automáticamente para las referencias.
%
% Archivos:
%   datos.tex          -> título, autores, ISBN, editorial...
%   preliminares/*.tex -> portada, página legal, resumen, abstract, prólogo...
%   capitulos/*.tex    -> un archivo por capítulo
%   referencias.bib    -> fuentes bibliográficas (APA automático)
%   estructura.tex     -> diseño y formato APA (normalmente no se toca)
%%%%%%%%%%%%%%%%%%%%%%%%%%%%%%%%%%%%%%%%%

\documentclass[12pt,fleqn]{book} % APA: 12 pt

% Color del libro: escribe azul o verde
\newcommand{\tema}{azul}

%%%%%%%%%%%%%%%%%%%%%%%%%%%%%%%%%%%%%%%%%
% Estructura del libro (diseño + normas APA 7.ª edición, en español)
%
% Basado en "The Legrand Orange Book", versión 2.1.1
% Autor original: Mathias Legrand, con modificaciones de Vel (LaTeXTemplates.com)
% Licencia: CC BY-NC-SA 3.0
%
% Cambios: idioma español, citas y referencias APA 7 (biblatex-apa),
% tablas y figuras en formato APA, márgenes de 2,54 cm, Times 12 pt,
% interlineado 1,5 y paleta azul o verde clara en lugar del naranja.
%
% Normalmente NO necesitas editar este archivo: el color se elige en main.tex
% y los datos del libro en datos.tex.
%%%%%%%%%%%%%%%%%%%%%%%%%%%%%%%%%%%%%%%%%

%----------------------------------------------------------------------------------------
%	PAQUETES BÁSICOS Y PÁGINA
%----------------------------------------------------------------------------------------

% APA: márgenes de 2,54 cm (1 pulgada) por todos los lados, papel carta
\usepackage[letterpaper,margin=2.54cm,headsep=10pt,headheight=16pt]{geometry}

\usepackage{graphicx} % Imágenes
\graphicspath{{imagenes/\tema/}{imagenes/}} % Primero busca en la carpeta del color elegido

\usepackage{lipsum} % Texto de relleno (puedes quitarlo cuando ya no lo uses)

\usepackage{tikz} % Dibujos de la portada y los capítulos
\usetikzlibrary{babel} % Evita conflictos entre TikZ y el español

% Idioma español: "Tabla" en vez de "Cuadro", sin atajos que rompan TikZ
\usepackage[english,spanish,es-tabla,es-noshorthands,es-noquoting,es-nodecimaldot]{babel}
\addto\captionsspanish{%
  \renewcommand{\contentsname}{Tabla de contenido}%
  \renewcommand{\listtablename}{Lista de tablas}%
  \renewcommand{\listfigurename}{Lista de figuras}%
}

\usepackage{enumitem} % Listas
\setlist{nolistsep}

\usepackage{booktabs} % Líneas horizontales de tablas (APA: sin líneas verticales)

\usepackage{zref-totpages} % Número total de páginas para la página legal

%----------------------------------------------------------------------------------------
%	PALETA DE COLORES (se elige con \tema en main.tex)
%----------------------------------------------------------------------------------------

\usepackage{xcolor}
\def\temaverde{verde}
\ifx\tema\temaverde
  \definecolor{principal}{RGB}{61,179,134} % Verde claro
\else
  \definecolor{principal}{RGB}{58,154,217} % Azul claro
\fi
% Versión más oscura del color, para textos (mejor contraste sobre blanco)
\colorlet{acento}{principal!70!black}

%----------------------------------------------------------------------------------------
%	FUENTES E INTERLINEADO (APA: Times New Roman 12 pt)
%----------------------------------------------------------------------------------------

\usepackage{avant} % Avant Garde para títulos
\usepackage{mathptmx} % Times para el texto y las matemáticas
\usepackage{microtype}
\usepackage[utf8]{inputenc}
\usepackage[T1]{fontenc}

\usepackage{setspace}
\onehalfspacing % Interlineado 1,5. Usa \doublespacing si te piden doble espacio

\setlength{\parindent}{1.27cm} % APA: sangría de 1,27 cm en la primera línea
% Texto justificado. Si te piden alineado a la izquierda (APA estricto),
% descomenta las dos líneas siguientes:
%\usepackage{ragged2e}
%\AtBeginDocument{\RaggedRight\setlength{\parindent}{1.27cm}}

%----------------------------------------------------------------------------------------
%	CITAS Y REFERENCIAS APA 7
%----------------------------------------------------------------------------------------

\usepackage[autostyle,spanish=mexican]{csquotes} % Comillas “ ” en lugar de « »
\usepackage[style=apa,backend=biber]{biblatex}
\addbibresource{referencias.bib}
% APA usa "&" entre autores dentro del paréntesis y en la lista de referencias.
% Si tu docente pide "y" en todos los casos, descomenta estas dos líneas:
%\DeclareDelimFormat[parencite]{finalnamedelim}{\addspace\bibstring{and}\space}
%\DeclareDelimFormat[bib,biblist]{finalnamedelim}{\ifthenelse{\value{listcount}>\maxprtauth}{}{\addspace\bibstring{and}\space}}

\usepackage{calc}
\usepackage{makeidx}
% Índice alfabético opcional: descomenta \makeindex y \printindex (al final de main.tex)
%\makeindex

%----------------------------------------------------------------------------------------
%	TABLAS Y FIGURAS EN FORMATO APA
%----------------------------------------------------------------------------------------

% APA: "Tabla 1" en negrita, en su propia línea; el título en cursiva debajo;
% todo encima de la tabla o figura. Numeración corrida (1, 2, 3...), no por capítulo.
\usepackage{caption}
\captionsetup{
  labelsep=newline,
  labelfont=bf,
  textfont=it,
  justification=raggedright,
  singlelinecheck=false,
  font={stretch=1},
  position=top,
  skip=4pt
}
\counterwithout{table}{chapter}
\counterwithout{figure}{chapter}

% Nota debajo de una tabla o figura: \nota{Elaboración propia con datos de ...}
\newcommand{\nota}[1]{\par\vspace{4pt}{\raggedright\small\textit{Nota.}\ #1\par}}

% APA: citas de 40 palabras o más, en bloque con sangría y sin comillas
\renewenvironment{quote}{\list{}{\leftmargin=1.27cm\rightmargin=0pt\topsep=6pt}\item\relax}{\endlist}
\newenvironment{citalarga}{\begin{quote}}{\end{quote}}

% Resumen / Abstract
\newcommand{\palabrasclave}[1]{\par\medskip\hspace*{1.27cm}\textit{Palabras clave}: #1\par}
\newcommand{\keywords}[1]{\par\medskip\hspace*{1.27cm}\textit{Keywords}: #1\par}

% Encabezados APA de nivel 4 y 5 (en línea con el párrafo, terminan en punto)
\newcommand{\nivelcuatro}[1]{\par\indent{\sffamily\bfseries #1.}\ }
\newcommand{\nivelcinco}[1]{\par\indent{\sffamily\bfseries\itshape #1.}\ }

%----------------------------------------------------------------------------------------
%	TABLA DE CONTENIDO
%----------------------------------------------------------------------------------------

\usepackage{titletoc}

\contentsmargin{0cm}
\setcounter{tocdepth}{2} % Muestra capítulos, secciones y subsecciones (1, 1.1, 1.1.1)

% Partes
\titlecontents{part}[0cm]
{\addvspace{20pt}\centering\large\bfseries}
{}
{}
{}

% Capítulos
\titlecontents{chapter}[1.25cm]
{\addvspace{12pt}\large\sffamily\bfseries}
{\color{principal}\contentslabel[\Large\thecontentslabel]{1.25cm}\color{acento}}
{\color{acento}}
{\color{principal}\normalsize\;\titlerule*[.5pc]{.}\;\thecontentspage}

% Secciones
\titlecontents{section}[1.25cm]
{\addvspace{3pt}\sffamily\bfseries}
{\contentslabel[\thecontentslabel]{1.25cm}}
{}
{\hfill\color{black}\thecontentspage}
[]

% Subsecciones
\titlecontents{subsection}[1.25cm]
{\addvspace{1pt}\sffamily\small}
{\contentslabel[\thecontentslabel]{1.25cm}}
{}
{\ \titlerule*[.5pc]{.}\;\thecontentspage}
[]

% Lista de figuras
\titlecontents{figure}[0em]
{\addvspace{-5pt}\sffamily}
{\figurename~\thecontentslabel.\hspace*{1em}}
{}
{\ \titlerule*[.5pc]{.}\;\thecontentspage}
[]

% Lista de tablas
\titlecontents{table}[0em]
{\addvspace{-5pt}\sffamily}
{\tablename~\thecontentslabel.\hspace*{1em}}
{}
{\ \titlerule*[.5pc]{.}\;\thecontentspage}
[]

%----------------------------------------------------------------------------------------
%	MINI TABLA DE CONTENIDO EN LAS PÁGINAS DE PARTE
%----------------------------------------------------------------------------------------

\titlecontents{lchapter}[0em]
{\addvspace{15pt}\large\sffamily\bfseries}
{\color{acento}\contentslabel[\Large\thecontentslabel]{1.25cm}\color{acento}}
{}
{\color{acento}\normalsize\sffamily\bfseries\;\titlerule*[.5pc]{.}\;\thecontentspage}

\titlecontents{lsection}[0em]
{\sffamily\small}
{\contentslabel[\thecontentslabel]{1.25cm}}
{}
{}

\titlecontents{lsubsection}[.5em]
{\normalfont\footnotesize\sffamily}
{}
{}
{}

%----------------------------------------------------------------------------------------
%	ENCABEZADOS DE PÁGINA
%----------------------------------------------------------------------------------------

\usepackage{fancyhdr}

\pagestyle{fancy}
\makeatletter
% En los preliminares (Resumen, Prólogo...) solo se muestra el título, sin "Capítulo 0"
\renewcommand{\chaptermark}[1]{\markboth{\sffamily\normalsize\bfseries\if@mainmatter\chaptername\ \thechapter.\ \fi#1}{}}
\makeatother
\renewcommand{\sectionmark}[1]{\markright{\sffamily\normalsize\thesection\hspace{5pt}#1}{}}
\fancyhf{} \fancyhead[LE,RO]{\sffamily\normalsize\thepage}
\fancyhead[LO]{\rightmark}
\fancyhead[RE]{\leftmark}
\renewcommand{\headrulewidth}{0.5pt}
\renewcommand{\headrule}{{\color{principal}\hrule width\headwidth height\headrulewidth\vskip-\headrulewidth}}
\renewcommand{\footrulewidth}{0pt}
\fancypagestyle{plain}{\fancyhead{}\renewcommand{\headrulewidth}{0pt}}

% Quita el encabezado de las páginas en blanco al final de los capítulos
\makeatletter
\renewcommand{\cleardoublepage}{
\clearpage\ifodd\c@page\else
\hbox{}
\vspace*{\fill}
\thispagestyle{empty}
\newpage
\fi}

%----------------------------------------------------------------------------------------
%	ESTILOS DE RECUADROS (definiciones, ejemplos, ejercicios...)
%----------------------------------------------------------------------------------------

\usepackage{amsmath,amsfonts,amssymb,amsthm}

\newtheorem{notation}{Notación}[chapter]

% Recuadros con encabezado de color
\newtheoremstyle{principalnumbox}
{0pt}{0pt}{\normalfont}{}
{\small\bf\sffamily\color{acento}}
{\;}{0.25em}
{\small\sffamily\color{acento}\thmname{#1}\nobreakspace\thmnumber{\@ifnotempty{#1}{}\@upn{#2}}%
\thmnote{\nobreakspace\the\thm@notefont\sffamily\bfseries\color{black}---\nobreakspace#3.}}
\renewcommand{\qedsymbol}{$\blacksquare$}

\newtheoremstyle{blacknumex}
{5pt}{5pt}{\normalfont}{}
{\small\bf\sffamily}
{\;}{0.25em}
{\small\sffamily{\tiny\ensuremath{\blacksquare}}\nobreakspace\thmname{#1}\nobreakspace\thmnumber{\@ifnotempty{#1}{}\@upn{#2}}%
\thmnote{\nobreakspace\the\thm@notefont\sffamily\bfseries---\nobreakspace#3.}}

\newtheoremstyle{blacknumbox}
{0pt}{0pt}{\normalfont}{}
{\small\bf\sffamily}
{\;}{0.25em}
{\small\sffamily\thmname{#1}\nobreakspace\thmnumber{\@ifnotempty{#1}{}\@upn{#2}}%
\thmnote{\nobreakspace\the\thm@notefont\sffamily\bfseries---\nobreakspace#3.}}

\newtheoremstyle{principalnum}
{5pt}{5pt}{\normalfont}{}
{\small\bf\sffamily\color{acento}}
{\;}{0.25em}
{\small\sffamily\color{acento}\thmname{#1}\nobreakspace\thmnumber{\@ifnotempty{#1}{}\@upn{#2}}%
\thmnote{\nobreakspace\the\thm@notefont\sffamily\bfseries\color{black}---\nobreakspace#3.}}
\makeatother

\newcounter{dummy}
\numberwithin{dummy}{section}
\theoremstyle{principalnumbox}
\newtheorem{theoremeT}[dummy]{Teorema}
\newtheorem{problem}{Problema}[chapter]
\newtheorem{exerciseT}{Ejercicio}[chapter]
\theoremstyle{blacknumex}
\newtheorem{exampleT}{Ejemplo}[chapter]
\theoremstyle{blacknumbox}
\newtheorem{vocabulary}{Vocabulario}[chapter]
\newtheorem{definitionT}{Definición}[section]
\newtheorem{corollaryT}[dummy]{Corolario}
\theoremstyle{principalnum}
\newtheorem{proposition}[dummy]{Proposición}

%----------------------------------------------------------------------------------------
%	RECUADROS DE COLOR
%----------------------------------------------------------------------------------------

\RequirePackage[framemethod=default]{mdframed}

% Teorema
\newmdenv[skipabove=7pt,skipbelow=7pt,backgroundcolor=black!5,linecolor=principal,
innerleftmargin=5pt,innerrightmargin=5pt,innertopmargin=5pt,innerbottommargin=5pt,
leftmargin=0cm,rightmargin=0cm]{tBox}

% Ejercicio
\newmdenv[skipabove=7pt,skipbelow=7pt,rightline=false,leftline=true,topline=false,bottomline=false,
backgroundcolor=principal!10,linecolor=principal,
innerleftmargin=5pt,innerrightmargin=5pt,innertopmargin=5pt,innerbottommargin=5pt,
leftmargin=0cm,rightmargin=0cm,linewidth=4pt]{eBox}

% Definición
\newmdenv[skipabove=7pt,skipbelow=7pt,rightline=false,leftline=true,topline=false,bottomline=false,
linecolor=principal,innerleftmargin=5pt,innerrightmargin=5pt,innertopmargin=0pt,innerbottommargin=0pt,
leftmargin=0cm,rightmargin=0cm,linewidth=4pt]{dBox}

% Corolario
\newmdenv[skipabove=7pt,skipbelow=7pt,rightline=false,leftline=true,topline=false,bottomline=false,
linecolor=gray,backgroundcolor=black!5,
innerleftmargin=5pt,innerrightmargin=5pt,innertopmargin=5pt,innerbottommargin=5pt,
leftmargin=0cm,rightmargin=0cm,linewidth=4pt]{cBox}

% Recuadro de presentación (texto destacado en los preliminares)
\newmdenv[skipabove=7pt,skipbelow=7pt,rightline=false,leftline=true,topline=false,bottomline=false,
backgroundcolor=principal!8,linecolor=principal,
innerleftmargin=12pt,innerrightmargin=12pt,innertopmargin=12pt,innerbottommargin=12pt,
leftmargin=0cm,rightmargin=0cm,linewidth=4pt]{destacado}

\newenvironment{theorem}{\begin{tBox}\begin{theoremeT}}{\end{theoremeT}\end{tBox}}
\newenvironment{exercise}{\begin{eBox}\begin{exerciseT}}{\hfill{\color{principal}\tiny\ensuremath{\blacksquare}}\end{exerciseT}\end{eBox}}
\newenvironment{definition}{\begin{dBox}\begin{definitionT}}{\end{definitionT}\end{dBox}}
\newenvironment{example}{\begin{exampleT}}{\hfill{\tiny\ensuremath{\blacksquare}}\end{exampleT}}
\newenvironment{corollary}{\begin{cBox}\begin{corollaryT}}{\end{corollaryT}\end{cBox}}

%----------------------------------------------------------------------------------------
%	NOTA AL MARGEN (círculo con "N")
%----------------------------------------------------------------------------------------

\newenvironment{remark}{\par\vspace{10pt}\small
\begin{list}{}{\leftmargin=35pt\rightmargin=25pt}\item\ignorespaces
\makebox[-2.5pt]{\begin{tikzpicture}[overlay]
\node[draw=principal!60,line width=1pt,circle,fill=principal!25,font=\sffamily\bfseries,inner sep=2pt,outer sep=0pt] at (-15pt,0pt){\textcolor{acento}{N}};\end{tikzpicture}}
\advance\baselineskip -1pt}{\end{list}\vskip5pt}

%----------------------------------------------------------------------------------------
%	TÍTULOS DE SECCIÓN (niveles APA 2 y 3, número en el margen)
%----------------------------------------------------------------------------------------

\makeatletter
\renewcommand{\@seccntformat}[1]{\llap{\textcolor{principal}{\csname the#1\endcsname}\hspace{1em}}}
% Nivel 2 APA: alineado a la izquierda, negrita
\renewcommand{\section}{\@startsection{section}{1}{\z@}
{-4ex \@plus -1ex \@minus -.4ex}
{1ex \@plus.2ex }
{\normalfont\large\sffamily\bfseries}}
% Nivel 3 APA: alineado a la izquierda, negrita cursiva
\renewcommand{\subsection}{\@startsection {subsection}{2}{\z@}
{-3ex \@plus -0.1ex \@minus -.4ex}
{0.5ex \@plus.2ex }
{\normalfont\sffamily\bfseries\itshape}}
\renewcommand{\subsubsection}{\@startsection {subsubsection}{3}{\z@}
{-2ex \@plus -0.1ex \@minus -.2ex}
{.2ex \@plus.2ex }
{\normalfont\small\sffamily\bfseries}}
\renewcommand\paragraph{\@startsection{paragraph}{4}{\z@}
{-2ex \@plus-.2ex \@minus .2ex}
{.1ex}
{\normalfont\small\sffamily\bfseries}}

%----------------------------------------------------------------------------------------
%	PÁGINAS DE PARTE (opcionales)
%----------------------------------------------------------------------------------------

\newcommand{\@mypartnumtocformat}[2]{%
\setlength\fboxsep{0pt}%
\noindent\colorbox{principal!20}{\strut\parbox[c][.7cm]{\ecart}{\color{acento}\Large\sffamily\bfseries\centering#1}}\hskip\esp\colorbox{principal!40}{\strut\parbox[c][.7cm]{\linewidth-\ecart-\esp}{\Large\sffamily\centering#2}}}%
\newcommand{\@myparttocformat}[1]{%
\setlength\fboxsep{0pt}%
\noindent\colorbox{principal!40}{\strut\parbox[c][.7cm]{\linewidth}{\Large\sffamily\centering#1}}}%
\newlength\esp
\setlength\esp{4pt}
\newlength\ecart
\setlength\ecart{1.2cm-\esp}
\newcommand{\thepartimage}{}%
\newcommand{\partimage}[1]{\renewcommand{\thepartimage}{#1}}%
\def\@part[#1]#2{%
\ifnum \c@secnumdepth >-2\relax%
\refstepcounter{part}%
\addcontentsline{toc}{part}{\texorpdfstring{\protect\@mypartnumtocformat{\thepart}{#1}}{\partname~\thepart\ ---\ #1}}
\else%
\addcontentsline{toc}{part}{\texorpdfstring{\protect\@myparttocformat{#1}}{#1}}%
\fi%
\startcontents%
\markboth{}{}%
{\thispagestyle{empty}%
\begin{tikzpicture}[remember picture,overlay]%
\node at (current page.north west){\begin{tikzpicture}[remember picture,overlay]%
\fill[principal!20](0cm,0cm) rectangle (\paperwidth,-\paperheight);
\node[anchor=north] at (4cm,-3.25cm){\color{principal!40}\fontsize{220}{100}\sffamily\bfseries\@Roman\c@part};
\node[anchor=south east] at (\paperwidth-1cm,-\paperheight+1cm){\parbox[t][][t]{8.5cm}{
\printcontents{l}{0}{\setcounter{tocdepth}{1}}%
}};
\node[anchor=north east] at (\paperwidth-1.5cm,-3.25cm){\parbox[t][][t]{15cm}{\strut\raggedleft\color{acento}\fontsize{30}{30}\sffamily\bfseries#2}};
\end{tikzpicture}};
\end{tikzpicture}}%
\@endpart}
\def\@spart#1{%
\startcontents%
\phantomsection
{\thispagestyle{empty}%
\begin{tikzpicture}[remember picture,overlay]%
\node at (current page.north west){\begin{tikzpicture}[remember picture,overlay]%
\fill[principal!20](0cm,0cm) rectangle (\paperwidth,-\paperheight);
\node[anchor=north east] at (\paperwidth-1.5cm,-3.25cm){\parbox[t][][t]{15cm}{\strut\raggedleft\color{acento}\fontsize{30}{30}\sffamily\bfseries#1}};
\end{tikzpicture}};
\end{tikzpicture}}
\addcontentsline{toc}{part}{\texorpdfstring{%
\setlength\fboxsep{0pt}%
\noindent\protect\colorbox{principal!40}{\strut\protect\parbox[c][.7cm]{\linewidth}{\Large\sffamily\protect\centering #1\quad\mbox{}}}}{#1}}%
\@endpart}
\def\@endpart{\vfil\newpage
\if@twoside
\if@openright
\null
\thispagestyle{empty}%
\newpage
\fi
\fi
\if@tempswa
\twocolumn
\fi}

%----------------------------------------------------------------------------------------
%	TÍTULOS DE CAPÍTULO (nivel 1 APA, con imagen de cabecera)
%----------------------------------------------------------------------------------------

\newif\ifusechapterimage
\usechapterimagetrue
\newcommand{\thechapterimage}{}%
\newcommand{\chapterimage}[1]{\ifusechapterimage\renewcommand{\thechapterimage}{#1}\fi}%
\def\@makechapterhead#1{%
{\parindent \z@ \raggedright \normalfont
\ifnum \c@secnumdepth >\m@ne
\if@mainmatter
\begin{tikzpicture}[remember picture,overlay]
\node at (current page.north west)
{\begin{tikzpicture}[remember picture,overlay]
\node[anchor=north west,inner sep=0pt] at (0,0) {\ifusechapterimage\includegraphics[width=\paperwidth]{\thechapterimage}\fi};
\draw[anchor=west] (\Gm@lmargin,-9cm) node [line width=2pt,rounded corners=15pt,draw=principal,fill=white,fill opacity=0.6,inner sep=15pt]{\strut\makebox[22cm]{}};
\draw[anchor=west] (\Gm@lmargin+.3cm,-9cm) node {\huge\sffamily\bfseries\color{black}\thechapter. #1\strut};
\end{tikzpicture}};
\end{tikzpicture}
\else
\begin{tikzpicture}[remember picture,overlay]
\node at (current page.north west)
{\begin{tikzpicture}[remember picture,overlay]
\node[anchor=north west,inner sep=0pt] at (0,0) {\ifusechapterimage\includegraphics[width=\paperwidth]{\thechapterimage}\fi};
\draw[anchor=west] (\Gm@lmargin,-9cm) node [line width=2pt,rounded corners=15pt,draw=principal,fill=white,fill opacity=0.6,inner sep=15pt]{\strut\makebox[22cm]{}};
\draw[anchor=west] (\Gm@lmargin+.3cm,-9cm) node {\huge\sffamily\bfseries\color{black}#1\strut};
\end{tikzpicture}};
\end{tikzpicture}
\fi\fi\par\vspace*{270\p@}}}

\def\@makeschapterhead#1{%
\begin{tikzpicture}[remember picture,overlay]
\node at (current page.north west)
{\begin{tikzpicture}[remember picture,overlay]
\node[anchor=north west,inner sep=0pt] at (0,0) {\ifusechapterimage\includegraphics[width=\paperwidth]{\thechapterimage}\fi};
\draw[anchor=west] (\Gm@lmargin,-9cm) node [line width=2pt,rounded corners=15pt,draw=principal,fill=white,fill opacity=0.6,inner sep=15pt]{\strut\makebox[22cm]{}};
\draw[anchor=west] (\Gm@lmargin+.3cm,-9cm) node {\huge\sffamily\bfseries\color{black}#1\strut};
\end{tikzpicture}};
\end{tikzpicture}
\par\vspace*{270\p@}}
\makeatother

%----------------------------------------------------------------------------------------
%	AUTORES (portadilla)
%----------------------------------------------------------------------------------------

% \autor{Nombre completo}{Afiliación, ORCID, ciudad, correo}
\newcommand{\autor}[2]{%
  \par\noindent{\sffamily\bfseries\large\color{acento}#1}\par
  {\small\begin{spacing}{1}\noindent #2\end{spacing}}\par\medskip}

%----------------------------------------------------------------------------------------
%	ENLACES DEL PDF
%----------------------------------------------------------------------------------------

\usepackage{hyperref}
\hypersetup{hidelinks,breaklinks=true,bookmarksopen=false}
\usepackage{bookmark}
\bookmarksetup{
open,
numbered,
addtohook={%
\ifnum\bookmarkget{level}=0 % capítulo
\bookmarksetup{bold}%
\fi
\ifnum\bookmarkget{level}=-1 % parte
\bookmarksetup{color=principal,bold}%
\fi
}
}

%----------------------------------------------------------------------------------------
%	DATOS DEL LIBRO
%	Llena estos campos una sola vez: se usan en la portada, la página legal
%	y las propiedades del PDF.
%----------------------------------------------------------------------------------------

\newcommand{\titulolibro}{Lo público, a la vista}
\newcommand{\subtitulolibro}{} % Déjalo vacío {} si no lleva
\newcommand{\autoresportada}{Jaider Bautista\\ Eliud Capera } % Nombres en la portada
\newcommand{\editorial}{Editorial XXXXXXX}
\newcommand{\isbn}{XXX-XXX-XXXXX-X-X}
\newcommand{\edicion}{Primera edición}
\newcommand{\fechaedicion}{XXXXX 2026} % Mes y año
\newcommand{\anio}{2026}
\newcommand{\ciudad}{Florencia (Colombia)}
\newcommand{\disenocubierta}{Copilot (IA) mediante \textit{prompt} de XXXXXXX XXXXXX}

% Referencia APA del propio libro (página legal, "Cómo citar")
\newcommand{\comocitar}{Apellido, N. (Ed.). (\anio). \textit{\titulolibro}. \editorial.}

\hypersetup{pdftitle={\titulolibro}}

\begin{document}

%----------------------------------------------------------------------------------------
%	PRELIMINARES (numeración romana)
%----------------------------------------------------------------------------------------

\frontmatter

%----------------------------------------------------------------------------------------
%	PORTADA
%	El título, subtítulo y autores salen de datos.tex.
%	Para usar otra cubierta, reemplaza imagenes/azul/background.png
%	(o la de verde) por tu imagen con el mismo nombre.
%----------------------------------------------------------------------------------------

\begingroup
\thispagestyle{empty}
\begin{tikzpicture}[remember picture,overlay]
\node[anchor=north west,inner sep=0pt] at (current page.north west) {\includegraphics[width=\paperwidth,height=\paperheight]{background}};
\node[anchor=north,fill=white,fill opacity=0.75,text opacity=1,inner sep=0.8cm,minimum width=\paperwidth] at ([yshift=-9cm]current page.north){%
\begin{minipage}{0.8\paperwidth}\centering\sffamily\linespread{1.1}\selectfont
{\color{principal}\rule{3cm}{3pt}}\par\vspace{14pt}
{\fontsize{38}{44}\selectfont\bfseries\color{acento}\titulolibro\par}
\vspace{14pt}
{\Large \subtitulolibro\par}
\vspace{22pt}
{\large\bfseries \begin{tabular}{@{}c@{}}\autoresportada\end{tabular}\par}
\end{minipage}};
\end{tikzpicture}
\vfill
\endgroup
\cleardoublepage

%----------------------------------------------------------------------------------------
%	PORTADILLA: título y autores con su afiliación
%	Copia el bloque \autor{...}{...} por cada autor.
%----------------------------------------------------------------------------------------

\thispagestyle{empty}
\begin{center}
{\Huge\sffamily\bfseries\color{acento}\titulolibro\par}
\vspace{6pt}
{\Large\sffamily \subtitulolibro\par}
{\color{principal}\rule{0.6\linewidth}{1.5pt}\par}
\end{center}

\vspace{1.5cm}
\noindent{\sffamily\bfseries\Large Autores:}\par\bigskip

\autor{Jaider Bautista Rodriguez}{Estudiante de Ingeniería de Sistemas, Universidad de la Amazonia, Facultad de Ingeniería, grupo de investigación o semillero al que pertenece, ORCID: 0000-0000-0000-0000, Florencia, Colombia, \href{mailto:jaid.bautista@udla.edu.co}{jaid.bautista@udla.edu.co}}
\autor{XXXX XXXXXX XXXXX XXXXX}{Estudiante de Ingeniería de Sistemas, Universidad de la Amazonia, Facultad de Ingeniería, grupo de investigación o semillero al que pertenece, ORCID: 0000-0000-0000-0000, Florencia, Colombia, \href{mailto:x.xxxxx@udla.edu.co}{x.xxxxx@udla.edu.co}}



\vfill
\noindent\textit{Libro de investigación producto de la articulación entre el programa de Ingeniería de Sistemas y el proyecto de investigación “La gestión pública y los propósitos fundamentales de la hacienda pública en los municipios del departamento del Caquetá”, de la Universidad de la Amazonia.}

\cleardoublepage

%----------------------------------------------------------------------------------------
%	PÁGINA LEGAL (créditos)
%	Los campos con \comando se llenan en datos.tex.
%----------------------------------------------------------------------------------------

\thispagestyle{empty}
~\vfill
\begin{spacing}{1.15}
\small
\noindent\textbf{\titulolibro}\\
\noindent Libro de investigación producto de la articulación del programa de Ingeniería de Sistemas de la Universidad de la Amazonia.\\
Incluye bibliografía.

\medskip
\noindent \copyright\ \editorial\\
\ztotpages\ páginas\\
ISBN (digital): \isbn\\
Número y año de edición: \edicion, \fechaedicion.\\
Diseño de cubierta: \disenocubierta.\\
Diagramación y terminación: \editorial.

\medskip
\noindent\textbf{Cómo citar:}\\
\hangindent=1.27cm\hangafter=1\comocitar

\medskip
\noindent \edicion, \ciudad, \anio.

\medskip
\noindent Prohibida la reproducción total o parcial de esta obra con fines comerciales. Su utilización se debe realizar con carácter académico, siempre que se cite la fuente.

\medskip
\noindent\textit{“El contenido de esta obra corresponde al derecho de expresión de los autores y no compromete el pensamiento de la \editorial, ni genera su responsabilidad frente a terceros. Los autores asumen la responsabilidad por los derechos de autor y conexos contenidos en la obra, así como por la eventual información sensible publicada en ella”.}

\medskip
\noindent Florencia, Caquetá, Colombia
\end{spacing}
\cleardoublepage

%----------------------------------------------------------------------------------------
%	PRESENTACIÓN BREVE (texto destacado antes de la tabla de contenido)
%----------------------------------------------------------------------------------------

\thispagestyle{empty}
~\vfill
\begin{destacado}
\noindent Este libro de investigación es el resultado de la sinergia entre docentes y estudiantes del programa de Ingeniería de Sistemas de la Universidad de la Amazonia. Su finalidad es investigativa y formativa: enseñar, desde lo básico hasta lo avanzado, cómo la inteligencia de negocios con Power BI permite analizar la percepción ciudadana sobre la gestión pública, la transparencia y el control social en los 16 municipios del departamento del Caquetá..
\end{destacado}
\vfill
\cleardoublepage


% Tabla de contenido y listas
\chapterimage{chapter_head_2}
\pagestyle{empty}
\tableofcontents
\cleardoublepage
\listoftables     % Borra esta línea si el libro no tiene tablas
\cleardoublepage
\listoffigures    % Borra esta línea si el libro no tiene figuras
\cleardoublepage
\pagestyle{fancy}

%----------------------------------------------------------------------------------------
%	RESUMEN
%	APA: un solo párrafo sin sangría, de 150 a 250 palabras.
%	Palabras clave: de 3 a 5, en minúscula, separadas por comas.
%----------------------------------------------------------------------------------------

\chapter{Resumen}

\noindent Escribe aquí el resumen del libro en un solo párrafo: el propósito, el método o enfoque, los resultados principales y las conclusiones. \lipsum[1][1-8]

\palabrasclave{palabra uno, palabra dos, palabra tres, palabra cuatro, palabra cinco}

%----------------------------------------------------------------------------------------
%	ABSTRACT (traducción del resumen al inglés)
%----------------------------------------------------------------------------------------

\chapter{Abstract}

\begin{otherlanguage}{english}
\noindent Write the English version of the abstract here, in a single paragraph. \lipsum[1][1-8]

\keywords{keyword one, keyword two, keyword three, keyword four, keyword five}
\end{otherlanguage}

%----------------------------------------------------------------------------------------
%	PRÓLOGO
%----------------------------------------------------------------------------------------

\chapter{Prólogo}

Texto del prólogo. Normalmente lo escribe una persona distinta a los autores y presenta la obra, su contexto y su aporte. \lipsum[2]

%----------------------------------------------------------------------------------------
%	AGRADECIMIENTOS
%----------------------------------------------------------------------------------------

\chapter{Agradecimientos}

Texto de agradecimientos a las personas, grupos de investigación, semilleros e instituciones que hicieron posible esta obra. \lipsum[3]


%----------------------------------------------------------------------------------------
%	CUERPO DEL LIBRO (numeración arábiga)
%----------------------------------------------------------------------------------------

\mainmatter

% Parte opcional: descomenta para agrupar capítulos en partes
%\part{Título de la parte}

\chapterimage{chapter_head_2}
%----------------------------------------------------------------------------------------
%	CAPÍTULO 1
%	Este capítulo trae ejemplos de todo lo que necesitas en APA (citas, tablas,
%	figuras, citas largas, recuadros). Úsalo como guía y reemplaza el contenido.
%----------------------------------------------------------------------------------------

\chapter{Generalidades}

Texto introductorio del capítulo. En APA cada párrafo empieza con sangría de 1,27 cm; para un párrafo nuevo basta con dejar una línea en blanco en el código. \lipsum[1]

%------------------------------------------------

\section{Cómo citar en APA}

Las citas se generan solas a partir de \texttt{referencias.bib}:

\begin{itemize}
\item Cita narrativa: \textcite{apa2020} establece las normas de esta edición.
\item Cita parentética: las normas se actualizaron en su séptima edición \parencite{apa2020}.
\item Con número de página: \parencite[p.~45]{apellido2023}.
\item Varias fuentes a la vez: \parencite{apellido2023,autor2024}.
\item Tres o más autores (sale \enquote{et al.} automáticamente): \textcite{autor2024}.
\item Institución como autor: \parencite{dane2024}.
\item Otros tipos de fuente (capítulo, tesis, web, ley): \parencite{capitulo2022,tesis2021,web2025,ley617}.
\end{itemize}

Para una cita textual corta (menos de 40 palabras) usa comillas: según \textcite[p.~12]{apellido2023}, \enquote{la cita textual va entre comillas y con página}.

Las citas de 40 palabras o más van en bloque, con sangría y sin comillas:

\begin{citalarga}
Este es un ejemplo de cita en bloque. APA indica que las citas textuales de 40 palabras o más se presentan en un párrafo independiente, con sangría de 1,27 cm desde el margen izquierdo, sin comillas, y con la cita entre paréntesis al final, después del punto. \parencite[p.~171]{apa2020}
\end{citalarga}

%------------------------------------------------

\section{Tablas en formato APA}

Las tablas se mencionan en el texto antes de aparecer, por ejemplo: la Tabla~\ref{tab:ejemplo} resume los datos. El número va en negrita, el título en cursiva, ambos encima de la tabla; la nota va debajo.

\begin{table}[htbp]
\caption{Título de la tabla en cursiva y con mayúscula inicial}
\label{tab:ejemplo}
\begin{tabular*}{\linewidth}{@{\extracolsep{\fill}} l c c}
\toprule
\textbf{Municipio} & \textbf{Indicador 1} & \textbf{Indicador 2}\\
\midrule
Municipio A & 0,32 & 56,2 \\
Municipio B & 0,15 & 91,0 \\
Municipio C & 0,09 & 29,6 \\
\bottomrule
\end{tabular*}
\nota{Elaboración propia con datos de \textcite{dane2024}.}
\end{table}

%------------------------------------------------

\section{Figuras en formato APA}

Las figuras (gráficos, imágenes, diagramas, capturas de pantalla) siguen el mismo formato, como en la Figura~\ref{fig:ejemplo}.

\begin{figure}[htbp]
\caption{Título de la figura en cursiva}
\label{fig:ejemplo}
\centering
\includegraphics[width=0.6\linewidth]{placeholder}
\nota{Tomado de \textcite{autor2024}. Si la figura es propia, escribe \enquote{Elaboración propia}.}
\end{figure}

%------------------------------------------------

\section{Subtítulos y listas}

\subsection{Subtítulo de nivel 3}

Los capítulos equivalen al nivel 1 de APA, las secciones (\verb|\section|) al nivel 2 y las subsecciones (\verb|\subsection|) al nivel 3.

\nivelcuatro{Encabezado de nivel 4} Si necesitas un nivel más, usa \verb|\nivelcuatro{...}|: va en negrita, con sangría y en la misma línea del párrafo.

\nivelcinco{Encabezado de nivel 5} Y \verb|\nivelcinco{...}| va en negrita cursiva.

\subsection{Listas}

\begin{enumerate}
\item Primer elemento.
\item Segundo elemento.
\item Tercer elemento.
\end{enumerate}

%------------------------------------------------

\section{Recuadros de la plantilla}

La plantilla conserva los recuadros de colores del diseño original, ahora en español.

\begin{definition}[Gestión pública]
Escribe aquí la definición de un concepto clave del libro \parencite{apellido2023}.
\end{definition}

\begin{example}[Nombre del ejemplo]
Texto de un ejemplo práctico.
\end{example}

\begin{exercise}
Una pregunta o actividad para el lector.
\end{exercise}

\begin{remark}
Una nota o aclaración al margen del texto principal.
\end{remark}


\chapterimage{chapter_head_1}
%----------------------------------------------------------------------------------------
%	CAPÍTULO 2
%	Cambia el título y agrega tantos capítulos como necesites: copia este archivo
%	como 03-..., 04-... y agrégalo en main.tex con \input{capitulos/nombre}.
%----------------------------------------------------------------------------------------

\chapter{Título del capítulo de desarrollo}

\lipsum[4]

\section{Primera sección}

\lipsum[5]

\section{Segunda sección}

\lipsum[6]

\subsection{Subsección}

\lipsum[7]


\chapterimage{chapter_head_2}
%----------------------------------------------------------------------------------------
%	CAPÍTULO: RESULTADOS
%----------------------------------------------------------------------------------------

\chapter{Resultados}

Presenta aquí los resultados (por ejemplo, el tablero o \textit{dashboard}, tablas y figuras con su interpretación). \lipsum[8]


\chapterimage{chapter_head_1}
%----------------------------------------------------------------------------------------
%	CAPÍTULO: CONCLUSIONES
%----------------------------------------------------------------------------------------

\chapter{Conclusiones}

Escribe las conclusiones: qué se logró, para qué sirve y qué queda por investigar. \lipsum[9]


%----------------------------------------------------------------------------------------
%	REFERENCIAS (APA 7: orden alfabético, sangría francesa)
%----------------------------------------------------------------------------------------

\chapterimage{chapter_head_2}
\chapter*{Referencias}
\addcontentsline{toc}{chapter}{Referencias}
\markboth{\sffamily\normalsize\bfseries Referencias}{\sffamily\normalsize Referencias}
\printbibliography[heading=none]

%----------------------------------------------------------------------------------------
%	APÉNDICES (opcionales; APA los llama Apéndice A, B...)
%----------------------------------------------------------------------------------------

%\appendix
%\chapterimage{chapter_head_1}
%\chapter{Instrumento de recolección de datos}
%Contenido del apéndice.

%----------------------------------------------------------------------------------------
%	ÍNDICE ALFABÉTICO (opcional; marca palabras en el texto con \index{palabra})
%----------------------------------------------------------------------------------------

%\cleardoublepage
%\phantomsection
%\setlength{\columnsep}{0.75cm}
%\addcontentsline{toc}{chapter}{Índice alfabético}
%\printindex

\end{document}
